\section{Q\&A：有限上下文、外部记忆与下一步}

主体讲座在 1:00 左右结束，后续问题集中在三个边界：多模态 token 怎样区分来源，有限 context 怎样扩展，以及 Karpathy 接下来研究什么。因为部分观众提问不可辨认，本节只保留能够从回答中可靠恢复的教学点，不猜测缺失问题。

\subsection{Context 不一定只靠无限加长}

Karpathy 提出一种当时的 speculative idea：保持 context length 有限，让模型学会输出特殊 scratchpad 指令，把要记住的内容写入外部 notebook，之后再查询。它把“脑内短期状态”与“外部可读写 memory”分开，预示 tool use 和 memory-augmented agent 的接口思路，但课堂没有给出实现或实验。

\teachervoice{他的类比是：人不会把所有事实都放在脑中，而会使用记事本。模型也可以学会发出 start/end scratchpad 标记，由外部系统保存、索引和重新注入内容。重要的是协议必须可监督、可寻址、可验证；仅仅给模型一个更长 prompt 不等于拥有可靠长期记忆。}

\subsection{Epistemic honesty：不知道就明确说不知道}

被问到 ChatGPT internals 时，Karpathy 直接说自己不知道；被问到 S4 时也没有假装熟悉。他最后回到正在做的 nanoGPT：把 minGPT 重写成更小、更现代、更容易复现 GPT 的代码，并对语言产品的持续迭代感兴趣。这种回答为讲义设定了同样标准：公开证据之外不补全系统细节。

\lecturefigure{slide-61-thank-you.jpg}{结束页：理解架构之后，“go forth and transform”。}{Stanford Online 官方视频 01:00:31--01:00:45。}

结束页几乎没有技术内容，却完成了课程动作：从一堂导论转向后续领域讲座。经过前文，``transform'' 不再只是换一个 backbone，而是完成四步设计：定义元素、注入位置/类型、规定 connectivity、选择 objective；然后验证 attention 路由、MLP 计算、优化稳定性与硬件成本是否共同成立。

\subsection{本章小结}

Q\&A 把架构边界转化为系统接口：多模态需要 type/position information，有限 context 可以结合外部 memory 与工具协议，未知内部实现必须保持未知。nanoGPT 则提供了继续学习的最短可执行路径：先复现小模型，再用实验检验关于 attention、mask、context 与 scaling 的直觉。

